\section{Capture without abolition}
\label{sec:13.7}
Halting, restoring, rewriting and replacing are done to the bearer, and the record shows each. Capture is done to the institutions that hold it, and no single act shows it: the institution keeps its name and its powers and stops using them. Transparency and whistleblower protection are the two defenses, and an institution can keep both and metabolize both. There's no threshold that detects this. So: measure the trend, not the level. Is the cost of registering dissent rising while the underlying risk stays flat? Is the required justification getting longer? Are there more approvers? Are the roles that gate exceptions migrating toward the roles that set the standards? Is the rate at which reviewers drawn to a cleared body are denied clearance rising while the pool's composition stays flat? Each of those is read by class of person the deployment acts on, not for the deployment as a whole. For the populations a state sorted before any movement arrived (\S\ref{sec:3.1}), the level is already high and a trend over the whole deployment underreads exactly where the baseline point applies.

An instrument for detecting recuperation is unusually easy to recuperate, so whoever measures the trend must not also gate exceptions, and the measurement has to be refusable. A written justification can't show which is happening. The same hundred words are a cost the dissenter chose to pay or a toll the gatekeeper set, and only the trend over time tells them apart. A regime built on tripwires is a regime built on metrics, and chapter~\ref{sec:3}'s second operation is the metric replacing the thing measured: a bound that fires on a number is met by moving the number. So no bound in this book stands alone. Each is published beside the trend that would show it being gamed, and the party that computes the trend gates nothing.
